\documentclass[border=6pt]{standalone}
\usepackage{tikz}
\usepackage{pgfplots}
\pgfplotsset{compat=1.18}
\usetikzlibrary{positioning,arrows.meta}
\usepackage{xeCJK}
\setCJKmainfont{KaiTi_GB2312}
\setCJKsansfont{SimHei}
\begin{document}
\begin{tikzpicture}[font=\footnotesize]
\node[draw,rounded corners=2pt,fill=gray!10,text width=6.0cm,align=center,inner sep=4pt] (src) at (3.1,5.65)
 {\textbf{同一份源码（liboqs 0.10.0）+ clang 18}\\[1pt]{\scriptsize 源码层均写为算术掩码（无显式 if-else）}};

\node[draw,rounded corners=2pt,fill=red!7,text width=6.0cm,align=left,inner sep=4pt] (x86) at (3.1,3.70)
 {\textbf{x86\_64 裸金属产物}\\[2pt] 编译为按位测试后条件跳转（\texttt{bt} + \texttt{jae}）\\ 秘密标记：512 次命中 / 1、8、15 处位置};

\node[draw,rounded corners=2pt,fill=green!8,text width=6.0cm,align=left,inner sep=4pt] (arm) at (3.1,1.75)
 {\textbf{aarch64 原生产物}\\[2pt] 编译为位域提取与条件选择（\texttt{sbfx}、\texttt{and}、\texttt{csel}）\\ 秘密标记：0 次命中 / 0 处位置（机制不复现）};

\draw[-{Stealth[length=2mm]},thick] (src.south) -- (x86.north);
\draw[-{Stealth[length=2mm]},thick] (src.west) -- ++(-0.32,0) |- (arm.west);

\node[draw,dashed,rounded corners=2pt,fill=yellow!8,text width=6.0cm,align=center,font=\scriptsize] at (3.1,0.45)
 {架构机制对照：12 格反汇编与秘密标记结果};

\begin{scope}[xshift=9.6cm]
\begin{axis}[
  at={(0,0.35)}, anchor=south west,
  width=7.5cm, height=6.1cm,
  axis background/.style={fill=white},
  xmode=log, xmin=1e3, xmax=2e12,
  ymin=0.35, ymax=5.85,
  ytick={1,...,5},
  yticklabels={{high (10 ms)},{mid (1 ms)},{low (100 $\mu$s)},{回环 (veth)},{本地直调}},
  y tick label style={font=\scriptsize},
  xlabel={分辨 32 ns 级效应所需的每类样本数（对数坐标）},
  xlabel style={font=\scriptsize},
  tick label style={font=\scriptsize},
  grid=major, grid style={gray!15},
  axis lines=left,
  xbar, bar width=6.5pt,
]
\addplot[fill=blue!45,draw=blue!65!black] coordinates {
 (97586175893,1) (1003153317,2) (64636555,3) (38417,4) (13098,5)
};
\draw[red!75!black,dashed,thick] (axis cs:2000,0.35) -- (axis cs:2000,5.65);
\draw[blue!75!black,dashed,thick] (axis cs:10000,0.35) -- (axis cs:10000,5.65);
\node[anchor=west,font=\tiny,color=red!75!black,fill=white,inner sep=1pt] at (axis cs:55000,4.72) {红虚线：网络档实测样本 (2,000)};
\node[anchor=west,font=\tiny,color=blue!75!black,fill=white,inner sep=1pt] at (axis cs:55000,4.08) {蓝虚线：本地直调样本 (10,000)};
\end{axis}
\end{scope}
\end{tikzpicture}
\end{document}
